\section{Introduction}

\subsection{Macro Level: The Rise of Egalitarian Discourse and the Treatment of Others}
Over the past two decades, cultural norms governing interpersonal relationships and the treatment of others have undergone an unprecedented transformation \citep{swales1990genre,grossmann2023ai}. Traditional hierarchical, transactional, or authoritarian relationship models have been widely superseded by ideals of absolute egalitarianism, mutual emotional validation, and strict parity in labor, respect, and mutual accommodation. Concepts such as ``equal partnership,'' ``fair share,'' ``emotional labor,'' and ``holding space'' have become deeply embedded in contemporary social consciousness. In theory, this widespread cultural internalization of parity represents a profound ethical advancement, offering individuals shared normative frameworks to govern how they treat others fairly and prevent exploitation.

However, classic scholarship in social psychology establishes that moral declarations rarely guarantee moral conduct. As \citet{batson1997unmasking,batson1999moral} demonstrated in foundational behavioral experiments, human moral behavior is frequently dominated by \textit{moral hypocrisy}: the powerful psychological motivation to appear fair, ethical, and principled in one's treatment of others, while covertly evading the behavioral and material sacrifices required by genuine fairness. While social monitoring and public accountability can constrain blatant hypocrisy in formal settings, private interpersonal interactions operate under radically different conditions. Behind closed doors—across romantic partners, domestic cohabitants, family members, and close peers—power dynamics are often asymmetric, emotionally charged, and shielded from external scrutiny. Consequently, the public vocabulary of equality can readily be co-opted into an interpersonal lever, completely divorced from actual treatment of others \citep{cerulli2026dark}.

\subsection{Meso Level: Real-World Interpersonal Exploitation, Treatment of Others, and the Methodological Trap of Self-Reports}
Within the private sphere of real-world relationships—spanning romantic couples, domestic cohabitants, and close familial or peer bonds—an alarming behavioral paradox frequently unfolds in the interpersonal treatment of others. One individual vocally and persistently demands strict parity, policing others for perceived imbalances in domestic chores, emotional attentiveness, or shared burdens. Yet, beneath this righteous advocacy for equality, this same individual habitually engages in severe, one-sided exploitation when treating those closest to them. Such extractive behaviors manifest as offloading burdensome responsibilities, demanding unilateral emotional accommodations, weaponizing progressive psychological terminology to gaslight others, and retaliating with stonewalling or emotional withdrawal whenever reciprocal consideration is requested \citep{bates2024virtuous,savejnarong2022interpersonal}.

Investigating this hypocrisy in the treatment of others has historically posed a formidable methodological dilemma for behavioral scientists. Traditional laboratory experiments and self-report psychometric surveys suffer from near-fatal \textit{social desirability bias} and pervasive self-deception. When surveyed, individuals overwhelmingly present themselves as compassionate, fair-minded actors who treat others with utmost equity; almost no one willingly acknowledges that they exploit others while hiding behind egalitarian slogans. 

To overcome this measurement trap, researchers have increasingly recognized the analytical power of naturalistic digital repositories. Millions of individuals routinely turn to anonymous online forums—such as Reddit's \textit{r/AmItheAsshole}, \textit{r/relationship\_advice}, \textit{r/marriage}, and \textit{r/family\_issues}—to unburden themselves, recount private interpersonal conflicts and their treatment of others in unvarnished detail, and seek third-party adjudication. Crucially, in this study, \textbf{social media is treated strictly as an unobtrusive observational data source (a naturalistic behavioral mirror), rather than the phenomenon itself}. These rich narrative disclosures capture real-world interpersonal histories, documented behavioral interactions, and raw interpersonal grievances regarding how individuals actually treat others that are completely inaccessible through conventional survey instruments.

\subsection{Micro Level: Problematization, Conceptual Grounding, and FINER Research Questions}
Following \citet{alvesson2011generating}, this study problematizes the pervasive assumption in relationship science and social ethics: that the frequent articulation of egalitarian relationship standards reflects prosocial commitment to fair treatment of others. We deconstruct this assumption from first principles, proposing that in conflicted or exploitative relationships, parity rhetoric is actively operationalized as an offensive lever of interpersonal dominance, self-serving attribution, and asymmetric treatment of others.

To elucidate this dynamic, we integrate four complementary theoretical perspectives:
\begin{enumerate}[leftmargin=*]
    \item \textbf{Moral Hypocrisy Theory \citep{batson1997unmasking,batson2006business}:} Positing that individuals seek to harvest the emotional and relational dividends of being seen as an egalitarian partner in treating others while covertly maximizing self-interest.
    \item \textbf{Self-Serving Attribution Theory \& Moral Licensing \citep{monin2001moral,merritt2010moral,krettenauer2022when}:} Explaining how vocal championing of fairness provides ``moral credentials'' that license selfish transgressions, leading individuals to attribute relationship successes to their own virtues while attributing conflicts and grievances to the alleged shortcomings of others.
    \item \textbf{Virtuous Victimhood \& The Dark Ego Vehicle Principle \citep{ok2021signaling,bates2024virtuous}:} Demonstrating that individuals with subclinical Dark Triad traits strategically broadcast personal victimhood, vulnerability, and past emotional trauma as an unassailable moral trump card to demand unreciprocated accommodations from others while silencing dissent.
    \item \textbf{Discursive Psychology \& Interpersonal Accounts \citep{edwards1992discursive,scott1968accounts,billig1996arguing}:} Establishing that language, moral claims, and conversational accounts in relationship disputes are not merely passive reports of internal cognition, but performative discursive actions deployed to construct moral accountability, manage blame, and manufacture legitimacy for extractive interpersonal behavior.
\end{enumerate}

To empirically test this model without self-report distortions, we analyze $N = 142,580$ naturalistic real-world relationship dispute narratives. We combine transformer-based Natural Language Processing (Sentence-BERT embeddings and fine-tuned RoBERTa classifiers) with a triangulated multi-model Large Language Model (LLM) ``cultural informant'' pipeline \citep{ziems2024large,argyle2023out,grossmann2024generative} to decode latent narrative hypocrisy.

Specifically, this study addresses three FINER research questions:
\begin{itemize}[leftmargin=*]
    \item \textbf{RQ1 (Parity Demands vs. Relational Extraction):} To what extent does high-frequency deployment of Egalitarian Vocabulary and Parity Demands (\textit{EVP}) correlate with Exploitative Interpersonal Relationship behaviors (\textit{EIR}) in real-world relationship narratives, and how does Self-Serving Attribution (\textit{SSA}) mediate this dynamic?
    \item \textbf{RQ2 (Victimhood as an Interpersonal Lever):} How does Virtuous Victimhood Signaling (\textit{VVS}) moderate the relationship between parity rhetoric and self-serving justifications, enabling actors to extract concessions while neutralizing social censure in their treatment of others?
    \item \textbf{RQ3 (AI Cultural Informant Decoding):} How effectively do multi-model Large Language Models, functioning as synthesized cultural informants trained on vast human relational discourse, decode the latent cognitive dissonance between proclaimed egalitarianism and enacted interpersonal exploitation?
\end{itemize}

The remainder of this manuscript is organized as follows. Section \ref{sec:lit_review} provides a thematic synthesis of recent (2022--2026) Q1 literature on interpersonal exploitation, attribution, and computational social science, presenting our Literature Synthesis Matrix and formal hypotheses. Section \ref{sec:methodology} describes the naturalistic corpus, construct operationalization, and AI informant protocol. Section \ref{sec:results} reports empirical PLS-SEM path estimations and dissonance classifications. Section \ref{sec:discussion} discusses theoretical contributions to relationship science, steelmanned rival explanations, severity-ranked limitations, clinical counseling implications, and future research trajectories.
